\section{Appendix}

This appendix provides the technical details of the main text. 

\subsection{Lumped Capacitance Model}

The TPS is divided into $N$ components with $M$ elements. For component $i$, define $\Omega_i$ as the spatial domain of the component, and $\partial\Omega_i$ as the boundary of the component. For each interface between two components, define $l(i)$ and $r(i)$ as the left and right components that form the interface, respectively. 

For interface $i$, define $\Gamma_i=\partial\Omega_{l(i)}\cap\partial\Omega_{r(i)}$,
\begin{subequations}
    \begin{align}
        \theta_i^{-} (t) &= \frac{1}{\left|\Gamma_i\right|}\int_{\Gamma_i} T^{-}_h(x,t)d\Gamma,\\
        \theta_i^{+}(t) &= \frac{1}{\left|\Gamma_i\right|}\int_{\Gamma_i} T^{+}_h(x,t)d\Gamma\\
        \Delta\theta_i(t) &= \theta_i^{-}(t) - \theta_i^{+}(t)
    \end{align}
\end{subequations}
Define the trace operators,
\begin{subequations}
    \begin{align}
        C_i^{-}\vu &= \frac{1}{\left|\Gamma_i\right|}\int_{\Gamma_i} \sum_{j=1}^{N_u} u_j(t)\phi_j(x)d\Gamma,\\
        C_i^{+}\vu &= \frac{1}{\left|\Gamma_i\right|}\int_{\Gamma_i} \sum_{j=1}^{N_u}u_j(t)\phi_j(x)d\Gamma
    \end{align}
\end{subequations}
Stacking all the $N_{\Gamma}$ interfaces, define the trace extraction operator as,
\begin{equation}
    \vPsi^+ = \begin{bmatrix}
        C_1^{-} - C_1^{+}\\
        \vdots\\
        C_{N_{\Gamma}}^{-} - C_{N_{\Gamma}}^{+}
    \end{bmatrix}\in\mathbb{R}^{N_{\Gamma}\times MP}
\end{equation}
where $M$ is the number of elements and $P$ is the number of basis functions per element. Therefore, it follows that,
\begin{equation}
    \vs(t) = \vPsi^+\vu(t)
\end{equation}
Thus,
\begin{equation}
    \vPi^+ = \begin{bmatrix}
        \vPhi^+ \\ 
        \vPsi^+
    \end{bmatrix}
\end{equation}

Introduce the lifting matrix,
\begin{equation}
    \vPsi\in\mathbb{R}^{MP\times N_{\Gamma}}
\end{equation}
such that,
\begin{equation}
    \vPhi^+\vPsi = \vzero,\quad \vPsi^+\vPsi = \vI
\end{equation}

Note that $\Phi\bvu$ contributes $E_{\Gamma}\bvu$ to the interface jump, therefore the state decomposition can be written as,
\begin{equation}
    \vu = \vPhi\bvu + \vPsi\left(\Deltavtheta - E_{\Gamma}\bvu\right) + \delta\vu = \tilde{\vPhi}\bvu + \vPsi\Deltavtheta + \delta\vu
\end{equation}
where $\tilde{\vPhi} = \vPhi - \vPsi E_{\Gamma}$, and $\delta\vu$ is the deviation from the average temperature in the components.

\subsection{Coarse-Graining Example}

Consider the three-layered TPS shown in Fig. [X] with three components $\left\{\Omega_j\right\}_{j=1}^3$ and three contact interfaces $\left\{\Gamma_k\right\}_{k=1}^3$. The set of elements for each component, and on the minus and plus sides of each interface are defined as,
\begin{subequations}
    \begin{align}
        \cV_j &= \left\{m\in\{1,\ldots,M\}:E_m\subseteq\Omega_j\right\},\\
        \cF_k^{-} &= \left\{m\in\{1,\ldots,M\}:E_m\cap\Gamma_k\neq\emptyset, E_m\subseteq\Omega_{l(k)}\right\},\\
        \cF_k^{+} &= \left\{m\in\{1,\ldots,M\}:E_m\cap\Gamma_k\neq\emptyset, E_m\subseteq\Omega_{r(k)}\right\}
    \end{align}
\end{subequations}
where $\Omega_{l(k)}$ and $\Omega_{r(k)}$ are the left and right components that form the interface $\Gamma_k$, respectively, with $\Gamma_k = \partial\Omega_{l(k)}\cap\partial\Omega_{r(k)}$.